\documentclass{article}
\usepackage{ctex}
\usepackage{amsmath, amsthm, amssymb}
\usepackage{hyperref}
\usepackage{graphicx}
\usepackage{bm}
\usepackage[thehwcnt = 1]{iidef}
\thecourseinstitute{自动化系}
\thecoursename{机器学习基础}
\theterm{2017年秋季学期}
\hwname{作业}
\slname{\heiti{解}}
\begin{document}
\courseheader
\name{赵丰}
\begin{enumerate}
\item
\begin{solution}
对于给定的$\epsilon,\delta$和模型$c$,以及分布$D$,产生了$m$个平面的样本点，样本点的标签若落在$c$内标1，$c$外标0。给出算法$\mathcal{A}$求出$c$的一个估计为包含所有标签为1的样本点的以原点为圆心的最小的圆。要证明对于$m>(1/\epsilon)\log(1/\delta)$有下式成立
\begin{equation}\label{eq:1}
\Pr_{S\sim D^m} [R(h_s) > \epsilon]<\delta
\end{equation}


首先可以知道$\mathcal{A}$产生的估计在样本点有限的情况下总是含在集合$c$内，$x$服从分布$D$,
如果$P(x\in c)<=\epsilon$,我们总假设$\epsilon$是小量，这种情况的含义是说我们取的分布$D$产生的点大部分在$c$外，是那些标签为0的点，而我们的算法只会把标签为1的点错判为标签为0的点，因为按照分布$D$产生的标签为1的点很少，所以根据代价函数的定义，有$R(h)<=\epsilon$（很小）恒成立，所以式\eqref{eq:1}左边恒为0，自然成立。

下面假设$P(x\in c)>\epsilon$,在这种情况下我们尝试估计
$\Pr_{S\sim D^m} [R(h_s) > \epsilon]$,因为$P(x\in c)>\epsilon$，所以$\exists 0<r_0<r$,使得$P(x | r_0 \leq ||x|| \leq r)=\epsilon$。下面假设我们的$\mathcal{A}$产生的圆的半径$r_1$介于$r_0$和$r$之间，那么所有判错的点出现的概率小于等于$P(x | r_1 \leq ||x|| \leq r)$小于等于$P(x | r_0 \leq ||x|| \leq r)=\epsilon$,这种情况同样得到\eqref{eq:1}左边为0。所以我们讨论最坏的情形，即$r_1$比$r_0$还小，在这种情况下事件$R(h_s) > \epsilon$可以推出任给m个样本点产生的$c$的估计的半径 $r_1<r_0$,我们计算一下后者的概率，也就是这$m$个样本点没有一个落在$r_0$和$r$之间就会产生这个错判，
对于单个$x$，$P(x|\,||x||<r_0 \text{或} ||x||>r)=1-\epsilon$,又因为$m$个样本是独立同分布的，所以估计的半径 $r_1<r_0$的概率为
$(1-\epsilon)^m$,根据两个随机事件的包含关系不难得到:
\begin{equation}
\Pr_{S\sim D^m} [R(h_s) > \epsilon] \leq  (1-\epsilon)^m
\end{equation}
 如果我们令$(1-\epsilon)^m<e^{-\epsilon m}<\delta$，即满足了$PAC$学习的定义，
从而推出$m>(1/\epsilon)\log(1/\delta)$
\end{solution}
\item
Gertrude 选圆的方法对$c$的估计不一定能保证圆心在$x_0$。
\begin{figure}[!ht]
\centering
\includegraphics[width=6cm]{P2.eps}
\caption{目标圆在$c$中的补不全落在$r_1\cup r_2 \cup r_3$示意图}\label{fig:1}
\end{figure}
Gertude 持有的命题的逆否命题为
如果 由样本点张成的目标圆与$r_1,r_2,r_3$都有交集，那么$P(x\in c)>\epsilon$小于$\epsilon$。
举一个反例：如图\ref{fig:1}所示，代价函数的取值为大圆面积减去小圆面积，差集不全含在$r_1\cup r_2 \cup r_3$内，而根据图中的面积比例差集的面积是有可能大于$r_1\cup r_2 \cup r_3$面积之和即$\epsilon$的。

因此Gertrude的方法有误。

\item
本题证明
$\forall \delta >0$,存在不快于多项式增长的函数$p(\frac{1}{\delta})$，使得当$m>p(\frac{1}{\delta})$时下式成立： 
\begin{equation}\label{eq:3modify}
\Pr_{S\sim D^m}(R(h_S)=0)>1-\delta
\end{equation}

设$Z=\{(x_1,y_1),\dots,(x_u,y_u)\}$，其中$u=|Z|$,设我们的分布$D$是一般的 Multinomial 分布 即满足 $P(X=x_i)=p_i$,但我们这里假设
$0<p_i<1$。因为若某一个$p_i$为零，即sample中不可能取到这个点，可以把它从$Z$中剔除，如果某一个$p_i=1$，则由概率归一化条件，
其他的$p_j$全为零，问题是平凡的。所以我们可以取到$0<p_m=\min_{1\leq i\leq u}\{p_i\}<1$,考虑到
\begin{equation}
R(h)=\Pr_{x \sim D}[h(x)\neq c(x)]=\sum_{i=1}^u 1_{h(x_i)\neq y_i}
\end{equation}
所以$R(h)$可能取到的最小值为$p_m$，且只能从有限个离散的值中取。
我们知道PAC 学习算法找到的$h_S$可以使得泛化误差尽可能小，即有
$\Pr_{S\sim D^m}(R(h_S)\leq \epsilon)\geq 1-\delta$,对于$m>p_1(\frac{1}{\epsilon},\frac{1}{\delta})$成立。
这里我们取小量$\epsilon =\frac{p_m}{2}$，则使得$\Pr_{S\sim D^m}(R(h_S)\leq \epsilon)$相当于
$\Pr_{S\sim D^m}(R(h_S)=0)$。
因此我们得到多项式$p=p_1(\frac{2}{p_m},\frac{1}{\delta})$，当$m>p$时我们有
\begin{equation*}
\Pr_{S\sim D^m}(R(h_S)=0)>1-\delta
\end{equation*}

\item
当$t=0$时结论显然成立;
若$t>0$,
记$Y=X-\mu+\lambda$,则$Y$均值为$\lambda$,方差为$\sigma^2$,

考虑$P(Y\geq t+\lambda)$,由Markov 不等式我们有:
\begin{align}
P(Y\geq t+\lambda) \leq & P(Y^2\geq (t+\lambda)^2)\\
\leq & \frac{E(Y^2)}{(t+\lambda)^2}\\
= & \frac{\sigma^2+\lambda^2}{(t+\lambda)^2}
\end{align}
取$\lambda=\frac{\sigma^2}{t}$，代入上式化简得到
\begin{equation}
P(X-\mu \geq t) = P(Y\geq t+\lambda) \leq \frac{\sigma^2}{t^2+\sigma^2}
\end{equation}
\end{enumerate}
\end{document}
